\newpage
\chapter{PENDAHULUAN} \label{Bab I}

\section{Latar Belakang} \label{I.Latar Belakang}

Kemajuan ilmu pengetahuan dan teknologi kini memberikan peluang yang signifikan untuk eksplorasi luar angkasa, baik secara global maupun di tingkat nasional. Di sejumlah negara maju, sistem pengamatan astronomi telah beralih ke otomatisasi total dengan menggunakan teleskop robotik yang dapat melacak objek di langit dengan akurasi tinggi. Sistem otomatis ini memainkan peran krusial dalam astronomi observasional saat ini, berkat kemampuannya untuk meningkatkan ketepatan dan memudahkan penunjukan objek di langit \cite{Hadi2022}. Tren ini tidak hanya membawa manfaat untuk penelitian profesional, tetapi juga menjadi semakin penting dalam dunia pendidikan dan komunitas astronomi amatir. Hal ini karena tren tersebut dapat memudahkan pengamat dalam memahami dinamika objek langit dengan cara yang lebih tepat dan jelas.  

Namun, teleskop otomatis yang tersedia di pasaran juga tidak lepas dari tantangan. Berdasarkan wawancara awal di Observatorium Astronomi ITERA Tabel~\ref{lampiran:wwc}, beberapa teleskop otomatis menghadapi kesulitan pada tahap konfigurasi awal, khususnya dalam hal penyelarasan (\textit{alignment}) dan kalibrasi yang bisa memakan waktu dan memerlukan keterampilan teknis. Jika teleskop mengalami pergeseran,maka proses konfigurasi harus diulang. Penelitian sebelumnya juga menyebutkan bahwa teleskop otomatis sering kali memiliki konfigurasi awal yang cukup sulit bagi pengguna dengan keterampilan terbatas \cite{Dyer2020} \cite{Rangana2022}.

Permasalahan tersebut menjadi penting karena keterbatasan pada sistem teleskop otomatis saat ini dapat membatasi partisipasi masyarakat dalam pengembangan ilmu astronomi. Apabila tidak segera dicari solusi, maka pemanfaatan teleskop di kalangan pelajar, peneliti muda, maupun komunitas amatir akan tetap kurang optimal. Padahal, perkembangan teknologi lokal memungkinkan terciptanya instrumen yang lebih adaptif dan mudah digunakan sehingga dapat meningkatkan literasi astronomi di Indonesia.

Sebagai solusi, penelitian ini menawarkan pengembangan teleskop robotik berbasis dudukan alt-azimuth dengan sistem kontrol otomatis melalui aplikasi Android sebagai antarmuka utama. Aplikasi ini berfungsi sebagai pusat kendali untuk mengatur pergerakan teleskop, memilih objek langit, dan memantau orientasi secara real-time melalui komunikasi websocket dengan mikrokontroler. Dudukan alt-azimuth dipilih karena mekanismenya sesuai dengan arah ketinggian (altitude) dan azimut (azimuth) \cite{Bradley2014} \cite{Roy2018}, sehingga lebih mudah diimplementasikan. Sistem dilengkapi sensor MPU9250 dan MPU6050 yang memadukan gyroscope dan akselerometer untuk memperoleh data orientasi, sehingga teleskop dapat menyesuaikan posisi otomatis saat dinyalakan maupun ketika terjadi pergeseran.

MPU9250 dan MPU6050 memiliki kelemahan berupa \textit{noise}, \textit{drift}, dan ketidakakuratan data \cite{Alfian2021}. Untuk mengatasinya, digunakan \textit{Mahony Filter}, yaitu filter berbasis pengendali PI (\textit{Proporsional-Integral}) dengan parameter gain proporsional (kp) dan integral (ki) \cite{Brndle2025}. Filter ini menggabungkan data giroskop dan akselerometer untuk menghasilkan estimasi orientasi yang lebih akurat. Prinsip kerjanya adalah menyelaraskan vektor gravitasi hasil akselerometer dengan estimasi orientasi yang dihitung \cite{Brndle2025} \cite{SivaSivamani2024}.

Kinerja teleskop robotik dievaluasi berdasarkan beberapa parameter operasional kunci. Hal ini mencakup akurasi \textit{pointing} (penunjukan), kemampuan \textit{tracking} (pelacakan objek secara kontinu), stabilitas terhadap gangguan, serta waktu respons antar target. Untuk kuantifikasi kesalahan, metrik \textit{Mean Absolute Error} (MAE) dan \textit{Root Mean Square Error} (RMSE) digunakan untuk mengukur selisih sudut rata-rata dan menilai ketahanan sistem terhadap lonjakan kesalahan \cite{Mayer2024}. \textit{Mahony Filter} dievaluasi untuk meningkatkan akurasi data dari sensor inersia. Evaluasi ini melibatkan perbandingan estimasi orientasi (sumbu \textit{roll}, \textit{pitch}, \textit{yaw}) sebelum dan sesudah penyaringan terhadap nilai referensi, dengan menggunakan metrik MAE dan RMSE. Selain itu, diuji pula pengaruh variasi parameter gain proporsional (kp) terhadap waktu konvergensi dan ketahanan filter terhadap akselerasi non-gravitasi. Pengujian aplikasi menggunakan \textit{unit testing} dan \textit{blackbox testing}.

Penelitian ini membahas perancangan teleskop robotik berbasis dudukan \textit{alt-azimuth} yang dilengkapi sensor IMU serta \textit{Mahony Filter}. Metode ini digunakan untuk meningkatkan akurasi pembacaan orientasi sudut, dengan evaluasi kinerja sistem dilakukan menggunakan metrik MAE (\textit{Mean Absolute Error}) dan RMSE (\textit{Root Mean Square Error}) untuk membandingkan hasil sebelum dan sesudah penerapan filter. Hasil yang diharapkan adalah teratasinya keterbatasan teleskop manual sekaligus terwujudnya instrumen astronomi otomatis berbasis teknologi lokal yang dapat dimanfaatkan secara luas.

\section{Rumusan Masalah} \label{I.Rumusan Masalah}
Berdasarkan latar belakang yang telah diuraikan, rumusan masalah dalam penelitian ini dapat dirumuskan sebagai berikut:

\begin{enumerate}[noitemsep]
	\item Bagaimana merancang dan membangun teleskop robotik dua sumbu berbasis \textit{Alt-Azimuth} yang dapat menyesuaikan posisi akibat pergeseran atau perpindahan, serta mendukung pelacakan objek antariksa?
		
	\item Bagaimana hasil evaluasi integrasi sensor IMU dengan penerapan 
	\textit{Mahony Filter} dalam meningkatkan akurasi estimasi sudut pada teleskop robotik?
\end{enumerate}



\section{Tujuan Penelitian} \label{I.Tujuan Penelitian}
Berdasarkan rumusan masalah yang telah dirumuskan, tujuan dari penelitian ini adalah sebagai berikut:

\begin{enumerate}[noitemsep]
	\item Merancang dan membangun teleskop robotik dua sumbu berbasis \textit{Alt-Azimuth} yang mampu menyesuaikan posisi saat terjadi pergeseran atau perpindahan, serta dapat melakukan pelacakan terhadap objek antariksa.
	
	\item Mengevaluasi kinerja integrasi sensor IMU dengan penerapan \textit{Mahony Filter} dalam meningkatkan akurasi estimasi sudut pada sistem teleskop robotik.
\end{enumerate}





\section{Batasan Masalah} \label{I.Batasan}
Agar penelitian ini fokus dan terarah, maka dibatasi pada hal-hal berikut:

\begin{enumerate}[noitemsep]
	\item Sistem navigasi teleskop menggunakan tata koordinat  \textit{Alt-Azimuth}.
	
	\item Estimasi sudut terbatas pada data dari sensor IMU (\textit{gyroscope} dan akselerometer); tidak menggunakan kamera atau citra visual.
	
	\item Katalog bintang yang digunakan berasal dari data terbuka \textit{stellarium}.
	
	\item Teleskop yang digunakan menggunakan lensa dengan panjang fokus 700mm dengan diameter 76 mm dan okuler 20mm.
\end{enumerate}



\section{Manfaat Penelitian} \label{I.Manfaat}
Penelitian ini diharapkan dapat memberikan manfaat sebagai berikut:
	\begin{enumerate}[noitemsep]
		\item Manfaat Akademis 
		\begin{itemize}
			\item Menjadi referensi ilmiah terkait penerapan \textit{Mahony Filter} dalam sistem estimasi sudut pada aplikasi robotik.
			
			\item Memberikan pemahaman praktis mengenai perancangan sistem robotik berbasis koordinat \textit{Alt-Azimuth}.
		\end{itemize}
	
	\item Manfaat Praktis
	\begin{itemize}
		\item Menyediakan desain teleskop robotik sederhana, murah, dan efektif untuk keperluan edukasi dan pengamatan astronomi.
		
		\item Menjadi dasar pengembangan lebih lanjut dalam sistem teleskop otomatis berbasis teknologi lokal yang dapat diakses oleh pelajar, guru, peneliti, dan komunitas astronomi amatir.
	\end{itemize}
\end{enumerate}





\section{Sistematika Penulisan} \label{I.Sistematika}
Sistematika penulisan berisi pembahasan apa yang akan ditulis disetiap Bab. Sistematika pada umumnya berupa paragraf yang setiap paragraf mencerminkan bahasan setiap Bab. \par

\subsection*{Bab I}
Bab ini berisikan penjelasan latar belakang dari topik penelitian yang berlangsung, rumusan masalah dari masalah yang dihadapi pada penjelasan di latar belakang, tujuan dari penelitian, batasan dari penelitian, manfaat dari hasil penelitian, dan sistematika penulisan tugas akhir. \par

\subsection*{Bab II}
Bab ini membahas mengenai tinjauan pustaka dari penelitan terdahulu dan dasar teori yang berkaitan dengan penelitian ini.

\subsection*{Bab III}
Bab ini berisikan penjelasan alur kerja sistem, alat dan data yang digunakan, metode yang digunakan, dan rancangan pengujian.